\subsection{Compatible square moments at a unit root}
\label{sec:unit-square-payments}

\begin{lemma}[A square with compatible moments]
\label{lem:unit-square-hermite}
Suppose $0\le H\le(c+X/2)^2$, $|X|\le U$, $c>U/2$, and
\[
 \E X\le-\alpha,\qquad \E X^2\le hU^2,\qquad
 0\le\alpha\le\sqrt h\,U,\quad U>0,\quad0\le h<1.
\]
Put $k=\alpha/U$ and
\[
 x_-=-U\frac{k+h}{1+k},\qquad
 \pi=\frac{h-k^2}{1+2k+h}.
\]
Then, for every integer $j\ge2$,
\begin{equation}\label{eq:unit-square-hermite:general}
 \E H^j\le
 (1-\pi)(c+x_-/2)^{2j}+\pi(c+U/2)^{2j}.
\end{equation}
In particular,
\begin{equation}\label{eq:unit-square-hermite}
 \E H^4\le
 (1-\pi)(c+x_-/2)^8+\pi(c+U/2)^8.
\end{equation}
If $U=0$, the bound for the $j$th moment is $c^{2j}$.
\end{lemma}
\begin{proof}
The quadratic tangent to $f(x)=(c+x/2)^{2j}$ at $x_-$ and agreeing
with $f$ at $U$ majorizes $f$ on $[-U,U]$. Its Hermite remainder
has the sign of $(x-x_-)^2(x-U)$, since $f'''\ge0$.
The quadratic coefficient is nonnegative by convexity, and the linear
coefficient is nonnegative because $x_-\le0$.
Hence its expectation is at most its value at the moment bounds
$\E X=-\alpha$, $\E X^2=hU^2$.
The two points $x_-,U$, with upper probability $\pi$, have exactly
these moments; $x_-\in[-U,0]$ and $0\le\pi\le1$.
This proves the general bound; $j=4$ gives the displayed eighth-power
formula. The zero case is direct.
\end{proof}

\begin{corollary}[The high-domain two-step path cost]
\label{cor:unit-square:two-step-path}
In the finite skew-kernel model, put \(p=1+K\one\) and
\(\bar f_2=T^2\one\). Then
\begin{equation}\label{eq:unit-square:two-step-path}
 \E(p\bar f_2)^8\le
 \begin{cases}
  35/8,&1/5\le d\le1/4,\\
  82/15,&1/4\le d\le1/3.
 \end{cases}
\end{equation}
\end{corollary}
\begin{proof}
Let \(r=K^2\one\). The nonnegative factors satisfy
\(p+\bar f_2=2+r\), so \(p\bar f_2\le(1+r/2)^2\).
The same actual \(r\) has
\(|r|\le\sqrt d\), \(\E r=-d\), and \(\E r^2\le\theta d\),
where \(\theta=13/32\). Apply
\eqref{eq:unit-square-hermite:general} with \(j=8\),
\(c=1,U=\sqrt d,\alpha=d,h=\theta\).
The negative atom contributes at most one, and
\(\pi=(\theta-d)/(1+2\sqrt d+\theta)\).
For \(d\ge1/4\), \(\pi\le5/77<1/15\) and
\((1+\sqrt d/2)^{16}<(13/10)^{16}<67\), giving \(82/15\).
For \(1/5\le d\le1/4\), use
\(\pi\le33/353<3/32\) and
\((1+\sqrt d/2)^{16}\le(5/4)^{16}<36\), giving \(35/8\).
The two last power comparisons follow from
\((13/10)^4<20/7\), \((20/7)^4<67\), and \((5/4)^8<6\).
\end{proof}

The following three scalar payments apply whenever the square in
Lemma~\ref{lem:unit-square-hermite} has the displayed moment data,
on $0\le d\le1/3$. Write $U=t\,u(d)$, $\alpha=t^2a(d)$,
and $c=C/D$, with $D=1$.
\Needspace{12\baselineskip}
\begingroup\small
\begin{longtable}{@{}clllll@{}}
\caption{Compatible square-moment data.}
\label{tab:unit-square-data}\\
\toprule
Row & $d$ & $u$ & $a$ & $c$ & $h$\\\midrule
I & $t^2$ & $1$ & $1/2$ & $1-3d/4$ & $13/32$\\
II & $63t^2/71$ & $1$ & $(63/71)(4/5-d)$ & $1-4d/5$ & $13/32$\\
III & $t^2$ & $1$ & $41/96$ & $1-37d/64$ & $3/10$\\\bottomrule
\end{longtable}
\endgroup
The targets are
\[
 \mathcal T_{\rm I}=(1-15d/32)^8,\quad
 \mathcal T_{\rm II}=(1-9d/20)^8,\quad
 \mathcal T_{\rm III}=(1-15d/32+3d^2/8)^8.
\]

Here is one polynomial certificate for all three targets. Define
\[
 \begin{split}
 P&=u+ta,\qquad Q=(1+h)u^2+2tau,\\
 L&=CP-\tfrac12Dtu(ta+hu),\qquad B=C+\tfrac12Dtu,\\
 N&=L^8+(hu^2-t^2a^2)B^8P^6,\qquad
 \mathcal D=D^8QP^6.
 \end{split}
\]
Substitution into \eqref{eq:unit-square-hermite} gives exactly
$N/\mathcal D$. The factors $P,Q,D$ are positive on each actual
interval. Set $\mathcal G=(\mathcal T\mathcal D-N)/t^2$, dividing out its exact endpoint factor.
Let $a_*$ be the left endpoint of the certificate interval and
$\widehat{\mathcal G}(t)=\mathcal G(t)/\mathcal G(a_*)$.
\Needspace{12\baselineskip}
\begingroup\small
\begin{longtable}{@{}cccc@{}}
\caption{Continuous polynomial certificates for Table~\ref{tab:unit-square-data}.}
\label{tab:unit-square-certificates}\\
\toprule
Row & Certificate interval for $t$ & Degree & Bernstein bound\\\midrule
I & $[0,577351/10^6]$ & $21$ & $B_j(\widehat{\mathcal G})\ge1/10$\\
II & $[0,612913/10^6]$ & $37$ & $B_j(\widehat{\mathcal G})\ge1/10$\\
III & $[0,577351/10^6]$ & $37$ & $B_j(\widehat{\mathcal G})\ge1/10$\\\bottomrule
\end{longtable}
\endgroup
In every row $\mathcal G(a_*)>1/2$. The coefficient formula is
\eqref{eq:bernstein-coefficients}, with one variable. Thus Lemma~\ref{lem:bernstein-certificate}
proves $\E H^4\le\mathcal T$ throughout the stated actual domains.
The rational certificate intervals include those domains, as is seen
by squaring their positive endpoints. At $d=0$, rows I--III have
$X=0,c=1$ and give equality.

